% Компилируется xelatex (рекомендуется) или pdflatex
\documentclass[12pt,a4paper]{article}

\usepackage{iftex}
\ifXeTeX
  \usepackage{fontspec}
  \IfFontExistsTF{Times New Roman}{\setmainfont{Times New Roman}}{%
    \IfFontExistsTF{Liberation Serif}{\setmainfont{Liberation Serif}}{\setmainfont{DejaVu Serif}}}
  \usepackage[russian]{babel}
\else
  \usepackage[T2A]{fontenc}
  \usepackage[utf8]{inputenc}
  \usepackage[russian]{babel}
\fi

\usepackage[left=25mm,right=15mm,top=20mm,bottom=20mm]{geometry}
\usepackage{longtable}
\usepackage{array}
\usepackage{booktabs}
\usepackage{graphicx}
\usepackage{float}
\usepackage{enumitem}
\setlist{nosep}
\setlength{\emergencystretch}{2em}
\setlist[description]{leftmargin=0pt,itemsep=3pt}

\newcolumntype{L}[1]{>{\raggedright\arraybackslash}p{#1}}

\newenvironment{reqtable}{%
  \begin{longtable}{L{14mm} L{146mm}}
  \toprule
  \textbf{ID} & \textbf{Требование} \\
  \midrule
  \endhead
  \bottomrule
  \endfoot
}{\end{longtable}}

\newenvironment{uctable}[1]{%
  \par\medskip\noindent\begin{minipage}{\linewidth}
  \textbf{#1}\par\smallskip
  \renewcommand{\arraystretch}{1.25}%
  \begin{tabular}{@{}L{38mm} L{124mm}@{}}
  \toprule
}{%
  \bottomrule
  \end{tabular}
  \end{minipage}\par\medskip
}

\title{Информационная система «Бдительность»\\[4pt]
\large Этап 1: описание предметной области и требования}
\author{}
\date{}

\begin{document}
\maketitle
\tableofcontents
\newpage

%=====================================================================
\section{Описание предметной области}

\subsection{Общая характеристика}

Предметная область основана на сеттинге игры \emph{Beholder}: тоталитарное
государство, в котором власть контролирует граждан через жилой фонд.
Рассматривается один многоквартирный дом. За дом отвечает управляющий,
назначенный государством. Он обслуживает жильцов и одновременно следит за
их благонадёжностью.

Отношения жильца с государством строятся на двух механизмах. Первый ---
доносы: любой жилец может сообщить о нарушении, допущенном соседом, а
управляющий проверяет сообщение и назначает наказание. Второй ---
государственные услуги: всё, что жилец получает от государства, оформляется
заявкой, и доступ к услугам зависит от благонадёжности жильца.

Благонадёжность выражается числом --- индексом лояльности. Нарушения
снижают его и закрывают доступ к услугам, а полезное для государства
поведение повышает его и открывает новые услуги. Таким образом, доносы и
услуги образуют единый цикл.

Система реализуется в виде веб-сайта, оформленного в стиле портала
государственных услуг.

\subsection{Участники}

\begin{description}
  \item[Жилец] --- гражданин, проживающий в квартире дома. Пополняет
    баланс, оплачивает квартплату, подаёт заявки на услуги и доносы на
    соседей.
  \item[Управляющий] --- должностное лицо, ответственное за дом. Заселяет
    и выселяет жильцов, рассматривает и исполняет заявки, рассматривает
    доносы и назначает наказания. В системе один управляющий.
\end{description}

\subsection{Ограничения}

\begin{enumerate}
  \item Система обслуживает один жилой дом.
  \item В системе две роли пользователей: жилец и управляющий. Управляющий один.
  \item Самостоятельная регистрация невозможна: учётную запись жилец создаёт
        только по коду приглашения.
  \item Перечень услуг, правила изменения ИЛ и статусы благонадёжности
        заданы в системе и не изменяются через интерфейс.
  \item Платежи принимаются только через ЮKassa в тестовом режиме.
  \item Сервер helios недоступен из сети Интернет, поэтому ЮKassa не может
        прислать уведомление об оплате: система сама периодически
        запрашивает статус платежа.
  \item В квартире проживает не более 10 жильцов.
\end{enumerate}

\subsection{Основные сущности и термины}

\begin{description}
  \item[Дом] --- объект учёта, состоящий из квартир. Система обслуживает
    один дом.
  \item[Квартира] --- жилое помещение, за которое начисляется квартплата.
    В квартире проживает до 10 жильцов; если жильцов меньше 10, в квартире
    есть свободное место.
  \item[Анкета] --- учётные данные жильца: фамилия, имя, отчество, дата
    рождения, телефон.
  \item[Код приглашения] --- одноразовый код, выдаваемый управляющим при
    заселении для активации учётной записи. Самостоятельная регистрация
    невозможна. Каждому жильцу выдаётся один код.
  \item[Индекс лояльности (ИЛ)] --- число от 0 до 100, отражающее
    благонадёжность жильца. Начальное значение --- 50. Изменяется
    автоматически по фиксированным правилам (раздел~\ref{sec:il}) и определяет
    статус благонадёжности, набор доступных услуг и очерёдность
    рассмотрения заявок.
  \item[Статус благонадёжности] --- категория жильца, определяемая
    значением ИЛ (раздел~\ref{sec:il}).
  \item[Коммунальный баланс] --- внутренний счёт жильца. Пополняется через
    платёжный сервис ЮKassa, расходуется на квартплату, пошлины и штрафы.
  \item[Задолженность] --- отрицательный коммунальный баланс. Возникает,
    если назначенный штраф больше остатка средств; погашается при
    следующем пополнении. Пошлина и квартплата при недостатке средств не
    списываются, поэтому задолженность образуют только штрафы.
  \item[Квартплата] --- ежемесячное начисление за квартиру в размере её
    тарифа. Начисление одно на квартиру: уведомление о нём получают все её
    жильцы, оплатить его может любой из них, и после оплаты остальным
    платить не нужно. Начисляется 1-го числа и оплачивается с баланса
    в течение 10 дней; просрочка оплаты является нарушением.
  \item[Услуга] --- предоставляемая государством услуга. Перечень услуг
    задан в системе. Для каждой услуги определены требуемый ИЛ, пошлина и
    форма результата.
  \item[Заявка] --- обращение жильца по услуге. Имеет статус, историю
    обработки и результат.
  \item[Результат заявки] --- итог исполнения: заключение управляющего,
    выданный документ или разрешение.
  \item[Донос] --- сообщение жильца о нарушении, допущенном другим жильцом.
    Может быть анонимным: автор сохраняется в системе, чтобы жилец видел
    свои доносы, но не раскрывается ни управляющему, ни нарушителю.
  \item[Наказание] --- решение управляющего о нарушении жильца:
    предупреждение, штраф или выселение. Назначается по одному или
    нескольким подтверждённым доносам на этого жильца либо без доноса; по
    одному доносу можно назначить несколько наказаний. Предупреждение и
    штраф снижают ИЛ нарушителя, выселение прекращает проживание.
    Назначенное наказание не изменяется; снять можно только предупреждение
    с помощью услуги «Снятие предупреждения из досье».
  \item[Досье] --- накопительное дело на жильца: доносы на него и поданные
    им, наказания, история ИЛ, заявки, платежи, заметки управляющего.
    Жильцу недоступно.
\end{description}

\subsection{Перечень услуг}

Перечень услуг фиксирован. Услуги упорядочены по возрастанию требуемого ИЛ:
чем ценнее услуга, тем выше должна быть благонадёжность жильца.

\begin{longtable}{L{11mm} L{50mm} L{9mm} L{20mm} L{55mm}}
\toprule
\textbf{ID} & \textbf{Услуга} & \textbf{ИЛ} & \textbf{Пошлина, руб.} & \textbf{Результат} \\
\midrule
\endhead
\bottomrule
\endfoot
SV01 & Бытовой ремонт & 0 & 0 & Отметка о выполнении работ \\
SV02 & Справка о проживании & 20 & 50 & Документ \\
SV03 & Продовольственный талон & 30 & 50 & Талон на указанный период \\
SV04 & Разрешение на приём гостей & 40 & 100 & Разрешение на указанные даты \\
SV05 & Разрешение на выход в комендантский час & 50 & 100 & Разрешение на указанные даты \\
SV06 & Замена мебели и бытовой техники & 60 & 300 & Отметка о выполнении работ \\
SV07 & Переселение в другую квартиру & 70 & 300 & Жилец переселён в выбранную квартиру со свободным местом \\
SV08 & Разрешение на выезд из города & 80 & 300 & Разрешение на период \\
SV09 & Снятие предупреждения из досье & 85 & 500 & Последнее действующее предупреждение снято, ИЛ повышен \\
\end{longtable}

При подаче заявки на талон или разрешение (SV03, SV04, SV05, SV08) жилец
указывает период. При подаче заявки на переселение (SV07) жилец выбирает
квартиру со свободным местом, отличную от своей. Заявку на SV09 можно подать, только
если у жильца есть действующее предупреждение.

\subsection{Индекс лояльности}\label{sec:il}

ИЛ изменяется только по правилам из таблицы ниже и всегда остаётся в
диапазоне от 0 до 100: изменение, выводящее ИЛ за границу, применяется
частично. Доносчик анонимного доноса ИЛ не получает и не теряет. Правило RL03
применяется к жильцу, оплатившему начисление, правило RL04 --- ко всем
жильцам квартиры, активировавшим учётную запись.

\begin{longtable}{L{14mm} L{110mm} L{24mm}}
\toprule
\textbf{ID} & \textbf{Событие} & \textbf{Изменение ИЛ} \\
\midrule
\endhead
\bottomrule
\endfoot
RL01 & Донос жильца подтверждён (кроме анонимного) & $+5$ \\
RL02 & Донос жильца признан ложным (кроме анонимного) & $-10$ \\
RL03 & Квартплата оплачена в течение 10 дней с даты начисления & $+1$ \\
RL04 & Квартплата не оплачена в течение 10 дней с даты начисления & $-5$ \\
RL05 & Жильцу вынесено предупреждение & $-10$ \\
RL06 & Жильцу назначен штраф & $-15$ \\
RL07 & Выполнена услуга SV09 «Снятие предупреждения из досье» & $+10$ \\
\end{longtable}

Статус благонадёжности определяется по текущему ИЛ и не хранится отдельно.

\begin{longtable}{L{40mm} L{40mm}}
\toprule
\textbf{ИЛ} & \textbf{Статус} \\
\midrule
\endhead
\bottomrule
\endfoot
80--100 & Образцовый \\
50--79 & Благонадёжен \\
20--49 & Под подозрением \\
0--19 & Неблагонадёжен \\
\end{longtable}

Критический уровень ИЛ --- 10: если ИЛ жильца опускается ниже 10,
управляющий получает уведомление о том, что жилец подлежит выселению.

\subsection{Модели состояний}

Статус заявки изменяется только по переходам из таблицы ниже.

\begin{longtable}{L{36mm} L{64mm} L{52mm}}
\toprule
\textbf{Статус заявки} & \textbf{Описание} & \textbf{Допустимые переходы} \\
\midrule
\endhead
\bottomrule
\endfoot
Подана & Заявка зарегистрирована, пошлина списана. & В работе, Отклонена, Отменена \\
В работе & Управляющий исполняет заявку. & Требует уточнения, Выполнена, Отклонена \\
Требует уточнения & Ожидается ответ жильца. & В работе \\
Выполнена & Результат оформлен. & --- \\
Отклонена & В услуге отказано с указанием причины. & --- \\
Отменена & Заявка отозвана жильцом, пошлина возвращена. & --- \\
\end{longtable}

Статус доноса изменяется только по переходам из таблицы ниже, поэтому
решение по рассмотренному доносу изменить нельзя.

\begin{longtable}{L{36mm} L{64mm} L{52mm}}
\toprule
\textbf{Статус доноса} & \textbf{Описание} & \textbf{Допустимые переходы} \\
\midrule
\endhead
\bottomrule
\endfoot
Подан & Донос зарегистрирован. & Подтверждён, Отклонён, Признан ложным \\
Подтверждён & Нарушение подтверждено. & --- \\
Отклонён & Оснований для наказания недостаточно. & --- \\
Признан ложным & Сведения доноса не соответствуют действительности. & --- \\
\end{longtable}

\subsection{Основные процессы}

\paragraph{Заселение.} Управляющий заселяет жильца в квартиру со свободным местом и
заполняет его анкету. Жилец получает начальный ИЛ 50, а управляющий --- код
приглашения, который передаёт жильцу. Жилец активирует учётную запись по
коду и получает доступ к личному кабинету.

\paragraph{Платежи.} Жилец пополняет баланс через ЮKassa. Каждый месяц
система начисляет квартплату, которую жилец оплачивает с баланса в течение
10 дней. Своевременная оплата повышает ИЛ, просрочка --- снижает.

\paragraph{Заявка на услугу.} Жилец выбирает доступную ему услугу и подаёт
заявку, оплачивая пошлину. Управляющий рассматривает заявку, при
необходимости запрашивает уточнение, затем исполняет её с оформлением
результата либо отклоняет с указанием причины. Жилец следит за статусом
заявки и получает результат в личном кабинете.

\paragraph{Донос.} Жилец подаёт донос на соседа. Управляющий рассматривает
его и либо подтверждает, либо отклоняет, либо признаёт ложным.
Подтверждённый донос повышает ИЛ неанонимного доносчика, ложный ---
снижает. По подтверждённым доносам управляющий назначает наказание; одно
наказание может опираться на несколько доносов. Управляющий также может
наказать жильца без доноса.

\paragraph{Выселение.} При падении ИЛ ниже 10 жилец подлежит выселению, и
управляющий получает уведомление. Управляющий выселяет жильца, его учётная запись
блокируется, место в квартире освобождается.

%=====================================================================
\section{Назначение информационной системы}

Без информационной системы обращения жильцов теряются, сроки их
рассмотрения не контролируются, доносы не проверяются, а решения
управляющего невозможно отследить. Система позволяет:

\begin{enumerate}
  \item вести учёт квартир и жильцов дома;
  \item предоставлять государственные услуги в электронном виде и
        отслеживать обработку заявок;
  \item ставить доступ к услугам в зависимость от благонадёжности жильца;
  \item принимать и рассматривать доносы, назначать наказания;
  \item автоматически оценивать благонадёжность жильцов через индекс
        лояльности;
  \item вести коммунальный баланс жильцов: пополнение, квартплату, пошлины
        и штрафы;
  \item накапливать сведения о жильцах в досье.
\end{enumerate}

%=====================================================================
\section{Функциональные требования}

Требования сгруппированы по акторам: жилец (U), управляющий (C), система (S).

\subsection{Жилец}

\subsubsection*{Учётная запись и профиль}
\begin{reqtable}
U001 & Жилец должен иметь возможность активировать учётную запись по коду приглашения, полученному от управляющего, и задать пароль. Код используется однократно. \\
U002 & Жилец должен иметь возможность войти в систему по логину и паролю и выйти из неё. Выселенный жилец войти в систему не может. \\
U003 & Жилец должен иметь возможность просматривать свою анкету и сведения о закреплённой за ним квартире. \\
U004 & Жилец должен иметь возможность просматривать текущее значение своего ИЛ, статус благонадёжности и историю изменений ИЛ с указанием даты и причины каждого изменения. Если изменение вызвано доносом, автор доноса не указывается. \\
\end{reqtable}

\subsubsection*{Баланс и платежи}
\begin{reqtable}
U005 & Жилец должен иметь возможность пополнить коммунальный баланс на выбранную сумму через ЮKassa. После оплаты жилец возвращается в личный кабинет и видит результат операции. \\
U006 & Жилец должен иметь возможность просматривать текущий баланс, сумму задолженности и историю операций: пополнения, квартплату, пошлины, штрафы и возвраты. \\
U007 & Жилец должен иметь возможность просматривать начисления квартплаты своей квартиры со статусом оплаты и плательщиком и оплачивать их с баланса. Оплатить начисление можно только при достаточном балансе; начисление, оплаченное другим жильцом квартиры, повторно не оплачивается. \\
\end{reqtable}

\subsubsection*{Заявки на услуги}
\begin{reqtable}
U008 & Жилец должен иметь возможность просматривать перечень услуг с описанием, требуемым ИЛ, размером пошлины и формой результата. Услуги, недоступные по текущему ИЛ, отображаются, но подать заявку на них нельзя. \\
U009 & Жилец должен иметь возможность подать заявку на доступную ему услугу, указав необходимые сведения: для талона и разрешений --- период, для переселения --- квартиру со свободным местом. При подаче с баланса списывается пошлина. Подать заявку при недостатке средств или наличии задолженности нельзя. \\
U010 & Жилец должен иметь возможность отменить свою заявку, пока управляющий не принял её в работу. При отмене пошлина возвращается на баланс. \\
U011 & Жилец должен иметь возможность просматривать список своих заявок и карточку каждой заявки: текущий статус, историю изменения статусов и комментарии управляющего. \\
U012 & Жилец должен иметь возможность ответить на запрос уточнения по заявке, после чего заявка возвращается управляющему на рассмотрение. \\
U013 & Жилец должен иметь возможность получить результат выполненной заявки: заключение управляющего и выданный документ или разрешение. При отклонении заявки жилец видит причину отказа. \\
U014 & Жилец должен иметь возможность просматривать полученные разрешения и документы, а также срок их действия. \\
\end{reqtable}

\subsubsection*{Доносы и наказания}
\begin{reqtable}
U015 & Жилец должен иметь возможность подать донос на другого жильца дома, указав суть нарушения, обстоятельства и, при желании, дату нарушения, которая не может быть позже текущей. Донос на самого себя подать нельзя. \\
U016 & Жилец должен иметь возможность подать донос анонимно. В этом случае система сохраняет автора, но его личность не раскрывается ни управляющему, ни нарушителю, а ИЛ доносчика по итогам рассмотрения не изменяется. \\
U017 & Жилец должен иметь возможность просматривать свои доносы и итог их рассмотрения: подтверждён, отклонён или признан ложным. \\
U018 & Жилец должен иметь возможность просматривать назначенные ему наказания с описанием нарушения и последствиями для баланса и ИЛ. Доносы, на основании которых назначено наказание, жильцу не показываются. \\
U019 & Жилец должен получать в личном кабинете уведомления об изменении статуса заявок, итогах рассмотрения доносов, пополнении баланса, начислении и просрочке квартплаты, наказаниях и изменениях ИЛ. Уведомления можно отметить прочитанными. \\
\end{reqtable}

\subsection{Управляющий}

\subsubsection*{Жилой фонд и жильцы}
\begin{reqtable}
C001 & Управляющий должен иметь возможность просматривать квартиры дома с указанием проживающих жильцов и их числа. \\
C002 & Управляющий должен иметь возможность заселить нового жильца в квартиру, где проживает меньше 10 жильцов, заполнив его анкету. При заселении жильцу назначается начальный ИЛ 50, а система выдаёт одноразовый код приглашения для активации учётной записи. \\
C003 & Управляющий должен иметь возможность просматривать список жильцов с их ИЛ, статусом благонадёжности и состоянием баланса. \\
C004 & Управляющий должен иметь возможность изменять анкетные данные жильцов. \\
C005 & Управляющий должен иметь возможность выселить жильца с указанием основания. Выселение фиксируется в досье как наказание, учётная запись жильца блокируется, а место в квартире освобождается. Выселенного жильца нельзя вернуть. \\
\end{reqtable}

\subsubsection*{Обработка заявок}
\begin{reqtable}
C006 & Управляющий должен иметь возможность просматривать очередь поступивших заявок. Заявки жильцов с более высоким ИЛ находятся выше в очереди. \\
C007 & Управляющий должен иметь возможность открыть заявку и просмотреть её содержание, сведения о заявителе и историю обработки. \\
C008 & Управляющий должен иметь возможность принять заявку в работу. После этого жилец не может её отменить. \\
C009 & Управляющий должен иметь возможность запросить у жильца уточнение по заявке с комментарием. \\
C010 & Управляющий должен иметь возможность отклонить заявку с обязательным указанием причины и решить, возвращается ли жильцу пошлина. \\
C011 & Управляющий должен иметь возможность завершить заявку, оформив результат: заключение и, если это предусмотрено услугой, документ или разрешение со сроком действия. \\
\end{reqtable}

\subsubsection*{Доносы и наказания}
\begin{reqtable}
C012 & Управляющий должен иметь возможность просматривать поступившие доносы. Автор анонимного доноса не отображается. \\
C013 & Управляющий должен иметь возможность подтвердить донос и назначить нарушителю наказание: предупреждение, штраф с указанием суммы или выселение. Наказание может назначаться по одному или нескольким подтверждённым доносам на этого жильца, а по одному доносу --- назначаться несколько наказаний. \\
C014 & Управляющий должен иметь возможность отклонить донос за недостатком оснований. В этом случае ИЛ участников не изменяется. \\
C015 & Управляющий должен иметь возможность признать донос ложным. В этом случае снижается ИЛ доносчика. \\
C016 & Управляющий должен иметь возможность самостоятельно зафиксировать нарушение жильца и назначить ему наказание без доноса. \\
C017 & Управляющий должен иметь возможность просматривать досье жильца: доносы на него и поданные им, назначенные наказания, историю ИЛ, заявки, платежи и заметки. Автор анонимного доноса в досье не отображается. \\
C018 & Управляющий должен иметь возможность добавлять в досье жильца собственные заметки. \\
\end{reqtable}

\subsection{Система}
\begin{reqtable}
S001 & Система должна ежемесячно начислять квартплату за каждую квартиру, где проживают жильцы, в размере тарифа квартиры со сроком оплаты 10 дней и уведомлять об этом всех жильцов квартиры. Повторное начисление за тот же месяц не выполняется. \\
S002 & Система должна фиксировать нарушение, если начисление не оплачено в течение 10 дней, и снижать по правилу RL04 ИЛ всех жильцов квартиры, активировавших учётную запись: до активации жилец не может оплатить начисление. \\
S003 & Система должна зачислять средства на баланс только после подтверждения оплаты от ЮKassa, ровно один раз на каждый платёж. \\
S004 & Система должна списывать назначенный штраф с баланса жильца. При недостатке средств баланс становится отрицательным (задолженность), и задолженность погашается при следующем пополнении. \\
S005 & Система должна автоматически изменять ИЛ по правилам RL01--RL07 при наказании, подтверждении или признании ложным доноса, своевременной оплате и просрочке квартплаты и выполнении услуги SV09, удерживая ИЛ в диапазоне от 0 до 100. \\
S006 & Система должна определять статус благонадёжности и набор доступных жильцу услуг по текущему значению ИЛ. \\
S007 & Система должна уведомлять управляющего о жильцах, чей ИЛ опустился ниже 10 и которые подлежат выселению. \\
S008 & Система должна вести историю всех операций с балансом, изменений ИЛ и изменений статусов заявок, которую нельзя изменить или удалить. \\
S009 & Система должна изменять статусы заявок и доносов только в соответствии с моделями состояний. Содержание поданного доноса, назначенное наказание и результат завершённой заявки не изменяются. \\
\end{reqtable}

%=====================================================================
\section{Нефункциональные требования}

\begin{reqtable}
N001 & \textbf{Платформа.} Веб-приложение на React, Spring и PostgreSQL, развёртываемое на сервере helios. \\
N002 & \textbf{Платежи.} Пополнение баланса выполняется через ЮKassa в тестовом режиме без использования реальных средств. \\
N003 & \textbf{Производительность.} Типовые операции выполняются без заметных для пользователя задержек при одновременной работе всех жильцов дома. \\
N004 & \textbf{Безопасность.} Пароли хранятся только в виде хеша bcrypt, права доступа проверяются на сервере, жилец получает доступ только к собственным данным, приложение защищено от распространённых веб-атак. \\
N005 & \textbf{Конфиденциальность.} Автор анонимного доноса не раскрывается. \\
N006 & \textbf{Целостность.} Финансовые операции, назначение наказаний и пересчёт ИЛ выполняются атомарно. \\
N007 & \textbf{Надёжность.} Ошибки обрабатываются корректно, пользователь получает понятные сообщения; повторная обработка подтверждения оплаты не приводит к двойному зачислению. \\
N008 & \textbf{Интерфейс.} Единый строгий стиль портала государственных услуг, интерфейс на русском языке. \\
N009 & \textbf{Совместимость.} Корректная работа в современных браузерах на компьютерах и мобильных устройствах. \\
N010 & \textbf{Сопровождаемость.} Многоуровневая архитектура с разделением представления, бизнес-логики и доступа к данным. \\
\end{reqtable}

%=====================================================================
\section{Прецеденты}

Ниже приведено описание прецедентов системы. Номер шага в альтернативном
потоке указывает на шаг основного потока, от которого он ответвляется.

\subsubsection*{Прецеденты актора «Жилец»}
\begin{uctable}{UC01. Активировать учётную запись}
Актор & Жилец \\
Цель & Получить доступ к личному кабинету. \\
Предусловия & Управляющий заселил жильца и выдал ему код приглашения. \\
Основной поток & 1.~Жилец открывает страницу активации. \newline 2.~Жилец вводит код приглашения. \newline 3.~Система проверяет код. \newline 4.~Жилец задаёт логин и пароль. \newline 5.~Система создаёт учётную запись, привязывает её к жильцу и открывает личный кабинет. \\
Альтернативные потоки & 3а. Код неверен или уже использован: система сообщает об ошибке, активация не выполняется. \\
Постусловия & Учётная запись активна, код приглашения больше не действует. \\
\end{uctable}

\begin{uctable}{UC02. Войти в личный кабинет}
Актор & Жилец \\
Цель & Начать работу с системой. \\
Предусловия & Учётная запись активирована. \\
Основной поток & 1.~Жилец вводит логин и пароль. \newline 2.~Система проверяет данные и статус жильца. \newline 3.~Система открывает личный кабинет. \\
Альтернативные потоки & 2а. Данные неверны: система сообщает об ошибке. \newline 2б. Жилец выселен: вход запрещён. \\
Постусловия & Жилец авторизован в системе. \\
\end{uctable}

\begin{uctable}{UC03. Просмотреть перечень услуг}
Актор & Жилец \\
Цель & Узнать, какие услуги доступны жильцу. \\
Предусловия & Жилец авторизован. \\
Основной поток & 1.~Жилец открывает раздел услуг. \newline 2.~Система показывает перечень услуг с требуемым ИЛ, пошлиной и формой результата. \newline 3.~Система отмечает услуги, недоступные по текущему ИЛ жильца. \\
Альтернативные потоки & — \\
Постусловия & — \\
\end{uctable}

\begin{uctable}{UC04. Пополнить баланс}
Актор & Жилец \\
Цель & Внести средства на коммунальный баланс. \\
Предусловия & Жилец авторизован. \\
Основной поток & 1.~Жилец указывает сумму пополнения. \newline 2.~Система создаёт платёж в ЮKassa и перенаправляет жильца на страницу оплаты. \newline 3.~Жилец оплачивает платёж. \newline 4.~Система периодически запрашивает статус платежа в ЮKassa и получает подтверждение оплаты. \newline 5.~Система однократно зачисляет средства на баланс; если есть задолженность, она погашается в первую очередь. \newline 6.~Жилец возвращается в личный кабинет и видит результат. \\
Альтернативные потоки & 3а. Жилец отменяет оплату: платёж отменяется, баланс не изменяется. \\
Постусловия & Баланс увеличен, операция отражена в истории. \\
\end{uctable}

\begin{uctable}{UC05. Просмотреть баланс и историю операций}
Актор & Жилец \\
Цель & Проконтролировать состояние счёта. \\
Предусловия & Жилец авторизован. \\
Основной поток & 1.~Жилец открывает раздел баланса. \newline 2.~Система показывает текущий баланс, задолженность и список операций. \\
Альтернативные потоки & — \\
Постусловия & — \\
\end{uctable}

\begin{uctable}{UC06. Оплатить квартплату}
Актор & Жилец \\
Цель & Погасить начисление за квартиру. \\
Предусловия & Система начислила квартплату. \\
Основной поток & 1.~Жилец открывает список начислений своей квартиры. \newline 2.~Жилец выбирает неоплаченное начисление и подтверждает оплату. \newline 3.~Система списывает сумму с баланса жильца и отмечает начисление оплаченным с указанием плательщика. \newline 4.~Если оплата выполнена в течение 10 дней с даты начисления, система повышает ИЛ плательщика по правилу RL03. \newline 5.~Остальные жильцы квартиры получают уведомление об оплате. \\
Альтернативные потоки & 2а. Начисление уже оплачено другим жильцом квартиры: оплата не требуется. \newline 3а. Недостаточно средств: оплата не выполняется, система предлагает пополнить баланс. \\
Постусловия & Начисление оплачено. \\
\end{uctable}

\begin{uctable}{UC07. Подать заявку на услугу}
Актор & Жилец \\
Цель & Получить государственную услугу. \\
Предусловия & Жилец авторизован, услуга доступна по его ИЛ. \\
Основной поток & 1.~Жилец выбирает услугу в перечне. \newline 2.~Жилец заполняет сведения по заявке; для талона и разрешений указывает период, для переселения --- квартиру со свободным местом. \newline 3.~Система проверяет ИЛ, баланс, отсутствие задолженности и сведения заявки. \newline 4.~Система списывает пошлину и регистрирует заявку в статусе «Подана». \newline 5.~Заявка попадает в очередь управляющего. \\
Альтернативные потоки & 3а. ИЛ ниже требуемого: подача невозможна. \newline 3б. Недостаточно средств или есть задолженность: система предлагает пополнить баланс. \newline 3в. Не указан период, в квартире для переселения нет свободного места или у жильца нет действующего предупреждения для SV09: подача невозможна. \\
Постусловия & Заявка зарегистрирована, пошлина списана. \\
\end{uctable}

\begin{uctable}{UC08. Отменить заявку}
Актор & Жилец \\
Цель & Отказаться от поданной заявки. \\
Предусловия & Заявка находится в статусе «Подана». \\
Основной поток & 1.~Жилец открывает заявку и выбирает отмену. \newline 2.~Система переводит заявку в статус «Отменена». \newline 3.~Система возвращает пошлину на баланс. \\
Альтернативные потоки & 1а. Заявка уже принята в работу: отмена недоступна. \\
Постусловия & Заявка отменена, пошлина возвращена. \\
\end{uctable}

\begin{uctable}{UC09. Ответить на запрос уточнения}
Актор & Жилец \\
Цель & Дополнить заявку по требованию управляющего. \\
Предусловия & Заявка находится в статусе «Требует уточнения». \\
Основной поток & 1.~Жилец получает уведомление и открывает заявку. \newline 2.~Жилец читает комментарий управляющего и вводит ответ. \newline 3.~Система возвращает заявку в статус «В работе». \\
Альтернативные потоки & — \\
Постусловия & Заявка снова на рассмотрении у управляющего. \\
\end{uctable}

\begin{uctable}{UC10. Получить результат заявки}
Актор & Жилец \\
Цель & Узнать итог рассмотрения заявки. \\
Предусловия & Заявка завершена или отклонена. \\
Основной поток & 1.~Жилец получает уведомление и открывает заявку. \newline 2.~Система показывает заключение управляющего и выданный документ или разрешение. \newline 3.~Жилец при необходимости скачивает документ. \\
Альтернативные потоки & 2а. Заявка отклонена: система показывает причину отказа и сведения о возврате пошлины. \\
Постусловия & — \\
\end{uctable}

\begin{uctable}{UC11. Подать донос}
Актор & Жилец \\
Цель & Сообщить о нарушении, допущенном соседом. \\
Предусловия & Жилец авторизован. \\
Основной поток & 1.~Жилец выбирает жильца, на которого подаёт донос. \newline 2.~Жилец описывает нарушение и обстоятельства и при желании указывает дату нарушения. \newline 3.~Жилец указывает, подаётся ли донос анонимно. \newline 4.~Система регистрирует донос и передаёт его управляющему. \newline 5.~Жилец может отслеживать статус доноса. \\
Альтернативные потоки & 1а. Жилец выбрал самого себя: система не допускает подачу. \newline 2а. Дата нарушения позже текущей: система не допускает подачу. \\
Постусловия & Донос зарегистрирован и ожидает рассмотрения. \\
\end{uctable}

\begin{uctable}{UC12. Просмотреть историю ИЛ}
Актор & Жилец \\
Цель & Узнать свой уровень благонадёжности и причины его изменения. \\
Предусловия & Жилец авторизован. \\
Основной поток & 1.~Жилец открывает раздел ИЛ. \newline 2.~Система показывает текущий ИЛ, статус благонадёжности и историю изменений с датами, величиной и правилом каждого изменения. \\
Альтернативные потоки & 2а. Изменение вызвано доносом: автор доноса не указывается. \\
Постусловия & — \\
\end{uctable}

\subsubsection*{Прецеденты актора «Управляющий»}
\begin{uctable}{UC13. Заселить жильца}
Актор & Управляющий \\
Цель & Зарегистрировать нового жильца дома. \\
Предусловия & В доме есть квартира со свободным местом. \\
Основной поток & 1.~Управляющий выбирает квартиру, где проживает меньше 10 жильцов. \newline 2.~Управляющий заполняет анкету жильца. \newline 3.~Система создаёт запись о жильце с начальным ИЛ 50 и одноразовый код приглашения. \newline 4.~Управляющий передаёт код жильцу. \\
Альтернативные потоки & 2а. Анкета заполнена не полностью: система требует заполнить обязательные поля. \\
Постусловия & Жилец проживает в квартире и может активировать учётную запись. \\
\end{uctable}

\begin{uctable}{UC14. Выселить жильца}
Актор & Управляющий \\
Цель & Освободить место в квартире и прекратить доступ жильца к системе. \\
Предусловия & Жилец проживает в доме. \\
Основной поток & 1.~Управляющий открывает карточку жильца и выбирает выселение. \newline 2.~Управляющий указывает основание. \newline 3.~Система фиксирует выселение в досье как наказание, блокирует учётную запись и освобождает место в квартире. \\
Альтернативные потоки & 1а. ИЛ жильца опустился ниже 10: система уведомляет управляющего, и управляющий принимает решение о выселении. \\
Постусловия & Жилец выселен, место в квартире свободно. \\
\end{uctable}

\begin{uctable}{UC15. Посмотреть досье}
Актор & Управляющий \\
Цель & Изучить сведения о жильце. \\
Предусловия & Жилец зарегистрирован в системе. \\
Основной поток & 1.~Управляющий выбирает жильца. \newline 2.~Система показывает досье: доносы на жильца и поданные им, наказания, историю ИЛ, заявки, платежи и заметки. Автор анонимного доноса не показывается. \\
Альтернативные потоки & — \\
Постусловия & — \\
\end{uctable}

\begin{uctable}{UC16. Добавить заметку в досье}
Актор & Управляющий \\
Цель & Зафиксировать собственные наблюдения о жильце. \\
Предусловия & Управляющий открыл досье жильца. \\
Основной поток & 1.~Управляющий вводит текст заметки. \newline 2.~Система сохраняет заметку с датой. \\
Альтернативные потоки & — \\
Постусловия & Заметка добавлена в досье. \\
\end{uctable}

\begin{uctable}{UC17. Посмотреть список заявок}
Актор & Управляющий \\
Цель & Выбрать заявку для обработки. \\
Предусловия & Управляющий авторизован. \\
Основной поток & 1.~Управляющий открывает очередь заявок. \newline 2.~Система показывает необработанные заявки, упорядоченные по убыванию ИЛ заявителей, а при равном ИЛ --- по дате подачи. \newline 3.~Управляющий открывает карточку заявки с содержанием, сведениями о заявителе и историей обработки. \\
Альтернативные потоки & — \\
Постусловия & — \\
\end{uctable}

\begin{uctable}{UC18. Принять заявку в работу}
Актор & Управляющий \\
Цель & Начать исполнение заявки. \\
Предусловия & Заявка в статусе «Подана». \\
Основной поток & 1.~Управляющий открывает заявку и выбирает «Принять в работу». \newline 2.~Система переводит заявку в статус «В работе» и уведомляет жильца. \\
Альтернативные потоки & — \\
Постусловия & Заявка в работе, жилец больше не может её отменить. \\
\end{uctable}

\begin{uctable}{UC19. Запросить уточнение}
Актор & Управляющий \\
Цель & Получить от жильца недостающие сведения. \\
Предусловия & Заявка в статусе «В работе». \\
Основной поток & 1.~Управляющий вводит комментарий с вопросом. \newline 2.~Система переводит заявку в статус «Требует уточнения» и уведомляет жильца. \\
Альтернативные потоки & — \\
Постусловия & Заявка ожидает ответа жильца. \\
\end{uctable}

\begin{uctable}{UC20. Отклонить заявку}
Актор & Управляющий \\
Цель & Отказать в предоставлении услуги. \\
Предусловия & Заявка в статусе «Подана» или «В работе». \\
Основной поток & 1.~Управляющий выбирает отклонение и указывает причину. \newline 2.~Управляющий решает, возвращается ли пошлина. \newline 3.~Система переводит заявку в статус «Отклонена», при необходимости возвращает пошлину и уведомляет жильца. \\
Альтернативные потоки & 1а. Причина не указана: система не позволяет отклонить заявку. \\
Постусловия & Заявка отклонена. \\
\end{uctable}

\begin{uctable}{UC21. Завершить заявку}
Актор & Управляющий \\
Цель & Выдать жильцу результат услуги. \\
Предусловия & Заявка в статусе «В работе». \\
Основной поток & 1.~Управляющий оформляет заключение. \newline 2.~Если услуга предусматривает документ или разрешение, управляющий прикладывает его и указывает срок действия. \newline 3.~Система переводит заявку в статус «Выполнена» и уведомляет жильца. \newline 4.~Если услуга меняет данные жильца, система применяет изменения: при переселении (SV07) меняет квартиру жильца, при снятии предупреждения (SV09) помечает последнее действующее предупреждение снятым и повышает ИЛ по правилу RL07. \\
Альтернативные потоки & 2а. Для талона или разрешения не указан срок действия: завершение невозможно. \newline 4а. В квартире для переселения не осталось свободного места: завершение невозможно. \\
Постусловия & Заявка выполнена, результат доступен жильцу. \\
\end{uctable}

\begin{uctable}{UC22. Рассмотреть донос}
Актор & Управляющий \\
Цель & Проверить сообщение о нарушении и принять решение. \\
Предусловия & Поступил донос. \\
Основной поток & 1.~Управляющий открывает донос и досье подозреваемого. \newline 2.~Управляющий оценивает обоснованность доноса. \newline 3.~Управляющий подтверждает донос. \newline 4.~Система переводит донос в статус «Подтверждён», повышает ИЛ неанонимного доносчика по правилу RL01 и уведомляет его. \newline 5.~Управляющий переходит к назначению наказания (UC23). \\
Альтернативные потоки & 3а. Оснований недостаточно: донос отклоняется, ИЛ участников не изменяется. \newline 3б. Донос ложный: донос признаётся ложным, ИЛ неанонимного доносчика снижается по правилу RL02. \\
Постусловия & Донос рассмотрен, доносчик видит итог. \\
\end{uctable}

\begin{uctable}{UC23. Назначить наказание}
Актор & Управляющий \\
Цель & Наказать жильца за нарушение. \\
Предусловия & Донос подтверждён, либо управляющий самостоятельно выявил нарушение. \\
Основной поток & 1.~Управляющий выбирает вид наказания: предупреждение, штраф или выселение. \newline 2.~Для штрафа управляющий указывает сумму. \newline 3.~Управляющий указывает подтверждённые доносы, на основании которых назначается наказание, либо назначает его без доноса. \newline 4.~Система фиксирует наказание в досье, снижает ИЛ нарушителя по правилу RL05 (предупреждение) или RL06 (штраф) и уведомляет его. \newline 5.~Для штрафа система списывает сумму с баланса. \\
Альтернативные потоки & 1а. Выбрано выселение: выполняется прецедент «Выселить жильца». \newline 3а. Донос не подтверждён или подан на другого жильца: он не может быть основанием наказания. \newline 5а. Недостаточно средств: баланс становится отрицательным, образуется задолженность. \\
Постусловия & Наказание назначено и отражено в досье. \\
\end{uctable}

%=====================================================================
\section{Используемые технологии}

Система строится по клиент-серверной архитектуре. Клиентская часть
представляет собой одностраничное веб-приложение, которое взаимодействует
с серверной частью через REST API. Серверная часть содержит бизнес-логику
и работает с базой данных.

\begin{longtable}{L{30mm} L{38mm} L{85mm}}
\toprule
\textbf{Компонент} & \textbf{Технология} & \textbf{Назначение} \\
\midrule
\endhead
\bottomrule
\endfoot
Клиентская часть & React & Пользовательский интерфейс личных кабинетов жильца и управляющего. \\
Серверная часть & Spring & Бизнес-логика (условия подачи заявок, правила изменения ИЛ, уведомления, последствия наказаний), REST API, аутентификация (проверка пароля по хешу bcrypt, завершение сеанса), разграничение доступа, периодический запрос статуса платежей в ЮKassa, ежемесячное начисление квартплаты и регистрация просрочек по расписанию. \\
База данных & PostgreSQL & Хранение данных о квартирах, жильцах, заявках, доносах, наказаниях, платежах и истории ИЛ; контроль целостности ограничениями и триггерами; функции для атомарных операций с деньгами и учётными записями и для сложных выборок. \\
Платёжный сервис & ЮKassa (тестовый режим) & Пополнение коммунального баланса. \\
\end{longtable}

\end{document}